\chapter{Metody Monte Carlo: symulacja, błąd i redukcja wariancji}
\label{ch:monte-carlo}
\paragraph{Powiązania i odsyłacze.} Prawo wielkich liczb i CTG pochodzą z tomu I (rozdziały~\ref{I-ch:lln} i~\ref{I-ch:clt}); symulacje procesów dyfuzyjnych omawia tom III, rozdział~\ref{III-ch:sde-properties}.


\section{Dlaczego losować, skoro mamy całkę?}
W statystyce często chcemy obliczyć całkę względem rozkładu a posteriori,
wartość predykcyjną albo ryzyko decyzji, a gęstość tego rozkładu jest
znana tylko z dokładnością do stałej normalizującej. Nawet gdy wartość
całki jest jednoznacznie zdefiniowana, obliczenie jej metodami
analitycznymi może być bardzo trudne. Metody Monte Carlo zastępują
całkowanie uśrednianiem odpowiednio dobranych zmiennych losowych.
W odróżnieniu od liczby danych $n$ z eksperymentu statystycznego,
liczbę \emph{symulacji} będziemy oznaczać przez $N$. Wzrost $N$ zmniejsza
błąd \emph{numeryczny}, lecz sam w sobie nie zwiększa informacji
statystycznej zawartej w pierwotnych danych.

\begin{definicja}[Podstawowy estymator Monte Carlo]
\label{def:mc-basic}
Niech $(\mathsf X,\mathcal A,P)$ będzie przestrzenią probabilistyczną,
$f\colon\mathsf X\to\R$ funkcją mierzalną i $f\in L^1(P)$.
Dla niezależnych $X_1,\ldots,X_N\sim P$ definiujemy
$\mu=\E_P[f(X_1)]=\int_{\mathsf X}f(x)\,\dd P(x)$ oraz
$\widehat\mu_N=N^{-1}\sum_{i=1}^N f(X_i)$.
Jeżeli $f\in L^2(P)$, oznaczamy
$\sigma_f^2=\Var_P(f(X_1))$.
\end{definicja}

\begin{twierdzenie}[Nieobciążoność, zgodność i centralne twierdzenie graniczne]
\label{thm:mc-lln-clt}
W warunkach definicji~\ref{def:mc-basic}, dla każdego $N\ge1$
$\E_P\widehat\mu_N=\mu$ oraz $\widehat\mu_N\to\mu$
prawie na pewno. Jeżeli ponadto $0<\sigma_f^2<\infty$, to
\[
 \Var_P(\widehat\mu_N)=\frac{\sigma_f^2}{N},\qquad
 \frac{\sqrt N(\widehat\mu_N-\mu)}{\sigma_f}
 \Rightarrow\Normal(0,1).
\]
Jeżeli $\sigma_f^2=0$, wówczas $\widehat\mu_N=\mu$ prawie na pewno.
\end{twierdzenie}
\begin{proof}
Liniowość wartości oczekiwanej daje nieobciążoność, a mocne prawo
wielkich liczb dla niezależnych zmiennych o jednakowym rozkładzie
-- zgodność. W przypadku $f\in L^2(P)$ niezależność zeruje
kowariancje między składnikami średniej, więc wariancja wynosi
$N^{-2}N\sigma_f^2=\sigma_f^2/N$.
Ostatnia zbieżność jest klasycznym centralnym twierdzeniem granicznym
z tomu I, zastosowanym do $f(X_i)$. Gdy wariancja jest zerowa,
$f(X_i)=\mu$ prawie na pewno dla każdego $i$; przeliczalność
rodziny pozwala uzyskać wspólne zdarzenie prawdopodobieństwa jeden.
\end{proof}

\begin{wniosek}[Przybliżony błąd Monte Carlo i liczba losowań]
\label{cor:mc-se}
W założeniach twierdzenia~\ref{thm:mc-lln-clt}, przy $N\ge2$, niech
$S_N^2=(N-1)^{-1}\sum_{i=1}^N(f(X_i)-\widehat\mu_N)^2$.
Wtedy $S_N^2\to\sigma_f^2$ prawie na pewno i
$\widehat{\mathrm{SE}}_{\rm MC}=S_N/\sqrt N$ jest zgodnym
oszacowaniem błędu standardowego Monte Carlo. Dla dużego $N$
przedział $\widehat\mu_N\pm z_{1-\alpha/2}S_N/\sqrt N$
ma asymptotyczne pokrycie $1-\alpha$ dla ustalonej całki $\mu$.
Aby zmniejszyć błąd standardowy dwukrotnie, przy niezmienionej
wariancji całkowanej funkcji trzeba około \emph{czterokrotnie}
zwiększyć $N$.
\end{wniosek}
\begin{proof}
Z mocnego prawa wielkich liczb dla $f(X_i)^2$ oraz dla $f(X_i)$
wynika $S_N^2\to\sigma_f^2$ prawie na pewno.
Twierdzenie~\ref{thm:mc-lln-clt} i twierdzenie Slutsky'ego
dają normalną granicę statystyki po podstawieniu $S_N$.
Zależność od $N^{-1/2}$ wyjaśnia ostatnie stwierdzenie.
\end{proof}

\begin{uwaga}[Błąd Monte Carlo nie jest błędem estymacji parametru]
Jeśli $\pi(\theta\mid x)$ jest posteriorem po danych $x$, to
całka $\mu(x)=\int_\Theta f(\theta)\pi(\theta\mid x)\,
\dd\lambda_d(\theta)$ jest \emph{ustaloną funkcją danych}.
Błąd przybliżenia $\widehat\mu_N(x)$ jest warunkowym błędem
numerycznym. Szerokość posterioru opisuje natomiast niepewność
$\theta$ przy danym $x$. Długi przebieg symulacji może bardzo dobrze
obliczyć szeroki posterior, lecz go nie zawęzi.
\end{uwaga}

\begin{przyklad}[Rzadkie zdarzenie -- ograniczenia naiwnego Monte Carlo]
\label{ex:mc-rare}
Jeśli $A$ ma prawdopodobieństwo $p$, bierzemy
$f=\ind{A}$; wtedy $\widehat p_N$ jest średnią z prób Bernoulliego,
a $\Var(\widehat p_N)=p(1-p)/N$. Względny błąd standardowy wynosi
$\sqrt{(1-p)/(Np)}$. Jeżeli $p=10^{-6}$ i $N=10^4$,
to oczekiwana liczba trafień to $0{,}01$, a prawdopodobieństwo
\emph{ani jednego trafienia} wynosi $(1-p)^N\approx0{,}99005$.
Otrzymane zero nie dowodzi, że zdarzenie jest niemożliwe.
\end{przyklad}

\section{Dwa sposoby zmniejszania kosztu: warunkowanie i zmienna kontrolna}
Idea warunkowania nawiązuje do twierdzenia Rao--Blackwella z rozdziału~\ref{ch:sufficiency}; tym razem służy redukcji błędu numerycznego, a nie ocenie ryzyka estymatora parametru.
\begin{twierdzenie}[Rao--Blackwell dla symulacji]
\label{thm:mc-rao-blackwell}
Niech $Z\in L^2(P)$ i $\mathcal G\subseteq\mathcal A$ będzie
pod-$\sigma$-ciałem. Wtedy
$\E[\E(Z\mid\mathcal G)]=\E Z$ oraz
$\Var(\E(Z\mid\mathcal G))\le\Var(Z)$,
z równością wtedy i tylko wtedy, gdy
$Z=\E(Z\mid\mathcal G)$ prawie na pewno.
Dlatego, jeżeli można obliczyć warunkową wartość oczekiwaną
bez kosztownej symulacji pozostałej losowości, warto symulować
$\E(Z\mid\mathcal G)$ zamiast $Z$.
\end{twierdzenie}
\begin{proof}
Pierwsza teza wynika z własności wieży, druga z rozkładu
$\Var(Z)=\E[\Var(Z\mid\mathcal G)]
 +\Var(\E[Z\mid\mathcal G])$.
Pierwszy składnik jest nieujemny i znika dokładnie wtedy,
gdy $Z=\E(Z\mid\mathcal G)$ prawie na pewno.
\end{proof}

\begin{twierdzenie}[Zmienna kontrolna]
\label{thm:control-variate}
Niech $Y,C\in L^2(P)$, $\E C=c$ jest dokładnie znane,
a $\Var(C)>0$. Dla ustalonej liczby $\beta\in\R$
zdefiniujmy $Z_\beta=Y-\beta(C-c)$.
Wtedy $\E Z_\beta=\E Y$ i
\[
 \Var(Z_\beta)=\Var(Y)-2\beta\Cov(Y,C)
                      +\beta^2\Var(C).
\]
Wariancję minimalizuje
$\beta_*=\Cov(Y,C)/\Var(C)$, a jej minimum wynosi
$\Var(Y)(1-\rho_{Y,C}^2)$, gdy $\Var(Y)>0$ i
$\rho_{Y,C}$ oznacza współczynnik korelacji.
\end{twierdzenie}
\begin{proof}
Odjęcie zmiennej o średniej zero nie zmienia wartości oczekiwanej.
Rozwinięcie wariancji kombinacji liniowej daje wielomian kwadratowy.
Po uzupełnieniu kwadratu
$\Var(Z_\beta)=\Var(C)(\beta-\beta_*)^2+
\Var(Y)-\Cov(Y,C)^2/\Var(C)$.
Wzór na korelację daje postać końcową.
\end{proof}

\begin{uwaga}[Pułapka doboru współczynnika na tych samych symulacjach]
Twierdzenie~\ref{thm:control-variate} dotyczy \emph{ustalonego}
$\beta$. Jeżeli wybieramy $\widehat\beta$ na podstawie tych
samych symulacji, które uśredniamy, nie wolno bez dowodu
przypisywać otrzymanej średniej dokładnej nieobciążoności.
Prosta metoda zachowania nieobciążoności polega na oszacowaniu
$\widehat\beta$ z niezależnej próby pilotażowej, a następnie
użyciu go jako stałej w nowej próbie.
\end{uwaga}

\begin{przyklad}[Średnia zmiennej lognormalnej]
Niech $Z\sim\Normal(0,1)$ i
$Y=e^Z$, $C=Z$. Wiemy, że $\E C=0$,
$\E Y=e^{1/2}$, $\Cov(Y,Z)=\E[Ze^Z]=e^{1/2}$
oraz $\Var(e^Z)=e^2-e$.
Najlepszy współczynnik kontrolny to $\beta_*=e^{1/2}$,
a wariancja maleje z $e^2-e$ do $e^2-2e$.
W przykładzie traktujemy $\E Y$ jako szukaną całkę,
mimo że znamy jej dokładną wartość i możemy sprawdzić błąd symulacji.
\end{przyklad}

\begin{stwierdzenie}[Zmienne antytetyczne]
\label{prop:antithetic}
Niech $U\sim\operatorname{Unif}(0,1)$, a
$f\in L^2((0,1),\lambda_1)$ będzie monotoniczna.
Wtedy $\Cov(f(U),f(1-U))\le0$.
Zatem estymator $[f(U)+f(1-U)]/2$ ma wariancję
nie większą niż $\Var(f(U))/2$.
Dla $N$ niezależnych par dzielimy tę wariancję przez $N$.
\end{stwierdzenie}
\begin{proof}
Niech $V$ będzie niezależną kopią $U$. Dla monotonicznej
$f$ funkcja $u\mapsto f(u)$ jest monotoniczna, a
$u\mapsto f(1-u)$ ma monotoniczność przeciwną.
Wobec tego
$[f(U)-f(V)][f(1-U)-f(1-V)]\le0$ prawie na pewno.
Po rozwinięciu i wykorzystaniu niezależności $U,V$
wartość oczekiwana lewej strony jest dwukrotnością
$\Cov(f(U),f(1-U))$, co daje tezę.
Wariancja średniej dwóch zmiennych o tej samej wariancji
wynosi $(\Var(f(U))+\Cov(f(U),f(1-U)))/2$.
\end{proof}

\section{Importance sampling: zmiana rozkładu losowania}
\begin{definicja}[Próbkowanie ważone]
\label{def:importance-sampling}
Niech $p,q$ będą gęstościami probabilistycznymi względem
tej samej miary $\sigma$-skończonej $\nu$ na $(\mathsf X,\mathcal A)$,
a $f\colon\mathsf X\to\R$ będzie mierzalna z
$\int_{\mathsf X}|f(x)|p(x)\,\dd\nu(x)<\infty$.
Zakładamy, że $q(x)>0$ dla $\nu$-prawie wszystkich $x$,
dla których $|f(x)|p(x)>0$. Definiujemy wagę
$w(x)=p(x)/q(x)$ dla $q(x)>0$ oraz $w(x)=0$ dla $q(x)=0$.
Dla niezależnych $X_i\sim q\,\nu$ estymatorem \emph{importance sampling}
jest $\widehat\mu_{N,\mathrm{IS}}=N^{-1}\sum_{i=1}^N f(X_i)w(X_i)$.
\end{definicja}

\begin{twierdzenie}[Nieobciążoność i wariancja importance sampling]
\label{thm:is-variance}
Przy założeniach definicji~\ref{def:importance-sampling} zachodzi
$\E_q\widehat\mu_{N,\mathrm{IS}}=\mu=
\int_{\mathsf X}f(x)p(x)\,\dd\nu(x)$.
Jeżeli ponadto
$\int_{\mathsf X}f(x)^2p(x)^2/q(x)\,\dd\nu(x)<\infty$
(z ilorazem rozumianym jako zero w punktach, w których $f(x)p(x)=0$), to
\[
 \Var_q(\widehat\mu_{N,\mathrm{IS}})=\frac1N
 \left(\int_{\mathsf X}\frac{f(x)^2p(x)^2}{q(x)}\,
       \dd\nu(x)-\mu^2\right).
\]
Spośród wszystkich gęstości $q$ dodatnich na $\{|f|p>0\}$
minimum drugiego momentu osiąga $q_*(x)=|f(x)|p(x)/A$,
gdzie $A=\int_{\mathsf X}|f(x)|p(x)\,\dd\nu(x)>0$.
Dla nieujemnego $f$ wariancja przy $q_*$ wynosi zero,
\emph{ale} konstrukcja $q_*$ wymaga znajomości całki $A=\mu$
i możliwości losowania z tego rozkładu.
\end{twierdzenie}
\begin{proof}
Z założenia nośnikowego $f(x)p(x)=0$ na $\{q=0\}$
prawie wszędzie. Zatem
$\E_q[fw]=\int_{\{q>0\}} f(x)p(x)\,\dd\nu(x)=\mu$.
Rachunek drugiego momentu i niezależność $X_i$ dają wariancję.
Dla dowolnej dopuszczalnej gęstości $q$, z nierówności
Cauchy'ego--Schwarza,
\[
 A^2=\left(\int_{\mathsf X}
       \frac{|f|p}{\sqrt q}\sqrt q\,\dd\nu(x)\right)^2
 \le\left(\int_{\mathsf X}\frac{f^2p^2}{q}\,\dd\nu(x)\right)
       \int_{\mathsf X}q(x)\,\dd\nu(x)
 =\int_{\mathsf X}\frac{f^2p^2}{q}\,\dd\nu(x).
\]
Równość zachodzi, gdy $|f|p/q$ jest stałe na nośniku $q$;
stąd $q_*=|f|p/A$. Gdy $f\ge0$, ważona wartość
$f(X)w(X)$ przy tym $q_*$ jest prawie na pewno równa $A$.
\end{proof}

\begin{uwaga}[Kiedy optymalny rozkład propozycji jest praktycznie nieosiągalny]
Gęstość $q_*$ z twierdzenia~\ref{thm:is-variance} jest optymalna dla
\emph{konkretnego} funkcjonału $f$, a nie jednocześnie dla wszystkich
całek względem $p$. Dla $f$ zmieniającej znak minimalna wariancja
wynosi $A^2-\mu^2$, która jest dodatnia, chyba że $f$ ma stały znak
$p$-prawie wszędzie. W praktycznych algorytmach przybliżamy $q_*$
rozwiązaniem parametrycznym, a potem oceniamy drugi moment wag:
duże ESS nie zastępuje weryfikacji pokrycia nośnika.
\end{uwaga}

\begin{uwaga}[Zero w mianowniku i nieistniejąca wariancja]
Może zachodzić $\E_q|fw|<\infty$, ale
$\E_q[(fw)^2]=\infty$. Wtedy estymator jest zgodny
na mocy mocnego prawa, lecz zwykły błąd standardowy
oparty na CLT nie ma uzasadnienia. W praktyce należy
kontrolować ogony wag i to, czy rozkład propozycji
$q$ jest dostatecznie szeroki na obszarach ważnych dla całki.
\end{uwaga}

\begin{twierdzenie}[Samonormalizowane importance sampling]
\label{thm:sn-is}
Niech $p$ będzie gęstością probabilistyczną,
$q>0$ na $\{p>0\}$ prawie wszędzie,
$f\in L^1(p\,\nu)$ oraz $X_i\iid q\,\nu$.
Zdefiniujmy $w_i=p(X_i)/q(X_i)$ oraz
$\widetilde\mu_N=\sum_{i=1}^Nw_i f(X_i)/\sum_{i=1}^Nw_i$
na zdarzeniu dodatniego mianownika (poza nim przyjmujemy
umownie wartość zero). Wtedy
$\widetilde\mu_N\to\mu$ prawie na pewno.
Jeżeli $\E_q[w^2(f-\mu)^2]<\infty$ i
$0<\Var_q(w(f-\mu))$, to
\[
 \sqrt N(\widetilde\mu_N-\mu)
 \Rightarrow\Normal\left(0,\E_q[w^2(f-\mu)^2]\right).
\]
Ten estymator nie jest na ogół dokładnie nieobciążony.
\end{twierdzenie}
\begin{proof}
Z warunku nośnikowego $\E_q w=1$ i
$\E_q(wf)=\mu$. MLLN dla obu średnich daje zgodność i
w szczególności dodatniość mianownika dla dużych $N$
prawie na pewno. Ponadto
\[
 \widetilde\mu_N-\mu=
 \frac{N^{-1}\sum_{i=1}^N w_i(f(X_i)-\mu)}
      {N^{-1}\sum_{i=1}^N w_i}.
\]
Licznik jest średnią niezależnych, wycentrowanych
zmiennych o skończonej wariancji wskazanej w tezie.
Z CLT i twierdzenia Slutsky'ego uzyskujemy normalną granicę.
Iloraz losowych średnich nie musi mieć wartości oczekiwanej
równej ilorazowi ich wartości oczekiwanych; już dla $N=1$
i różnych $p,q$ samonormalizowana średnia równa się
$f(X_1)$ przy $w_1>0$, więc jej wartość oczekiwana
może być inna niż $\mu$.
\end{proof}

\begin{uwaga}[Efektywna liczebność próby ważonej]
W praktyce oblicza się
$\mathrm{ESS}_{\rm w}=(\sum_iw_i)^2/\sum_iw_i^2\in[1,N]$
na zdarzeniu, że suma wag jest dodatnia.
Jest to \emph{wskaźnik koncentracji wag}, a nie dokładna
gwarancja błędu $\widetilde\mu_N$; przy wagach zerowych
liczymy tylko dodatnie składniki. Mała wartość oznacza,
że kilka losowań dominuje wynik. Sama duża wartość ESS
nie dowodzi, że rozkład propozycji pokrywa wszystkie
istotne obszary ani że istniała wariancja wag.
\end{uwaga}

\begin{przyklad}[Ekspozycyjna zmiana miary dla rzadkiego ogona]
Niech $Z\sim\Normal(0,1)$, $a>0$, a
$p=\Prob(Z>a)=\int_{(a,\infty)}\varphi(z)\,\dd\lambda_1(z)$,
gdzie $\varphi$ oznacza standardową gęstość normalną.
W naiwnej symulacji liczymy $\ind{\{Z>a\}}$.
W importance sampling losujemy $Y\sim\Normal(a,1)$ i liczymy
\[
 W(Y)=\ind{\{Y>a\}}\exp(-aY+a^2/2).
\]
Ponieważ $\varphi(y)/\varphi(y-a)=\exp(-ay+a^2/2)$,
$\E W(Y)=p$. Wybór przesunięcia $a$ sprawia, że
przekroczenie progu występuje pod propozycją
z prawdopodobieństwem $1/2$, zamiast bardzo rzadko.
Dobre ustawienie propozycji trzeba jednak ocenić przez
wariancję wag, a nie samą częstość przekroczeń.
\end{przyklad}

\section{Co kiedy stosować? Porównanie w jednym miejscu}
\begin{center}\small
\begin{tabular}{>{\raggedright\arraybackslash}p{.25\textwidth}>{\raggedright\arraybackslash}p{.30\textwidth}>{\raggedright\arraybackslash}p{.29\textwidth}}
\toprule
\textbf{Zadanie} & \textbf{Metoda do rozważenia} & \textbf{Kontrola jakości}\\
\midrule
Znamy rozkład, umiemy losować niezależnie & Zwykłe Monte Carlo & Błąd $S_N/\sqrt N$ przy skończonej wariancji\\
Rzadkie zdarzenie lub dominujące ogony & Importance sampling & Nośnik propozycji, ogony wag i ESS\\
Znamy średnią skorelowanej zmiennej & Zmienna kontrolna & Zysk wariancyjny; osobna próba pilotażowa\\
Umiemy scałkować część losowości & Warunkowanie Rao--Blackwella & Koszt wyliczenia warunkowej średniej\\
Monotoniczny funkcjonał jednego kwantyla & Losowania antytetyczne & Faktyczna kowariancja w parach\\
Posterior znany tylko do normalizacji & MCMC z rozdziału~\ref{ch:mcmc} & Diagnostyka wielu łańcuchów, autokorelacja\\
\bottomrule
\end{tabular}
\end{center}

\section{Laboratorium: błąd i rzadkie zdarzenia}
\begin{lstlisting}[language=Python,caption={Kompletny program: Monte Carlo, zmienna kontrolna i importance sampling dla ogona normalnego (Python, NumPy i SciPy).}]
import numpy as np
from scipy.stats import norm

rng = np.random.default_rng(2026)
n = 100_000
z = rng.standard_normal(n)
y = np.exp(z)
naive = y.mean()
control = (y - np.exp(0.5) * z).mean()
print("E[exp(Z)] exact:", np.exp(0.5))
print("naive estimate:", naive, "SE:", y.std(ddof=1) / np.sqrt(n))
print("control estimate:", control)

a = 4.0
x = rng.standard_normal(n)
rare_plain = (x > a).astype(float)
y_is = rng.normal(loc=a, scale=1.0, size=n)
w_is = (y_is > a) * np.exp(-a * y_is + a*a/2)
print("tail exact:", norm.sf(a))
print("plain:", rare_plain.mean(), "hits:", int(rare_plain.sum()))
print("importance sampling:", w_is.mean())
print("IS estimated SE:", w_is.std(ddof=1) / np.sqrt(n))

# Self-normalized importance sampling has generally nonzero bias.
weights = np.exp(-a * y_is + a*a/2)
snis = np.sum(weights * (y_is > a)) / weights.sum()
ess = weights.sum()**2 / np.sum(weights**2)
print("SNIS estimate:", snis, "weight ESS:", ess)
\end{lstlisting}

\begin{lstlisting}[language=Julia,caption={Kompletny program: redukcja wariancji przez zmienną kontrolną i estymacja rzadkiego ogona (Julia, Random i Statistics).}]
using Random, Statistics
Random.seed!(2026)
N = 100_000
z = randn(N)
y = exp.(z)
beta = exp(0.5)
println("exact E[exp(Z)]: ", beta)
println("plain: ", mean(y))
println("control: ", mean(y .- beta .* z))

a = 4.0
x = randn(N)
plain = mean(x .> a)
yis = randn(N) .+ a
wis = (yis .> a) .* exp.(-a .* yis .+ a^2 / 2)
println("plain tail estimate: ", plain)
println("importance sampling tail estimate: ", mean(wis))
println("IS estimated standard error: ", std(wis) / sqrt(N))
\end{lstlisting}

\begin{zadanie}[Dwie metody dla tej samej całki]
Niech $U\sim\operatorname{Unif}(0,1)$ i $f(u)=u^2$.
Porównać wariancje pojedynczej wartości $f(U)$ i
średniej antytetycznej $[f(U)+f(1-U)]/2$.
\par\emph{Rozwiązanie.} Mamy $\E U^2=1/3$,
$\E U^4=1/5$, więc $\Var(U^2)=4/45$.
Ponieważ $\E[U^2(1-U)^2]=1/30$,
$\Cov(U^2,(1-U)^2)=1/30-1/9=-7/90$.
Wariancja średniej pary jest
$\tfrac12(4/45-7/90)=1/180$.
Przy dwóch niezależnych próbach wariancja ich średniej
wynosiłaby $2/45$. Para antytetyczna jest więc przy
tej samej liczbie dwóch ewaluacji znacznie dokładniejsza.
\end{zadanie}


